\section{Hasil dan Pembahasan}

\subsection{Dampak Pemeliharaan Struktur}

Tabel~\ref{tab:hasil-chunker-paper} merangkum perbandingan \chunker{} pada cutoff 30. Pada \textit{User Manual}, \textit{structure-aware} dengan hybrid indexer memperoleh recall 0,740, lebih tinggi 0,055 daripada sliding-window dan 0,115 daripada recursive. Pada korpus hukum, \textit{legal structure-aware} dengan dense indexer memperoleh recall 0,6854, lebih tinggi 0,0509 daripada sliding-window dan 0,0124 daripada recursive.

\begin{table*}[t]
\centering
\caption{Hasil perbandingan \chunker{} pada cutoff $K=30$}
\label{tab:hasil-chunker-paper}
\scriptsize
\begin{tabular}{@{}llrrrr@{}}
\toprule
\textbf{Korpus} & \textbf{Strategi} & \textbf{Recall} & \textbf{Precision} & \textbf{NDCG} & \textbf{MRR}\\
\midrule
User Manual & Structure-aware hybrid & \textbf{0,740} & \textbf{0,048} & \textbf{0,333} & 0,231\\
User Manual & Recursive & 0,625 & 0,032 & 0,291 & 0,216\\
User Manual & Sliding-window & 0,685 & 0,041 & 0,322 & \textbf{0,232}\\
Peraturan & Legal structure-aware dense & \textbf{0,6854} & \textbf{0,1520} & 0,5040 & \textbf{0,5240}\\
Peraturan & Recursive & 0,6730 & 0,1453 & 0,4909 & 0,4937\\
Peraturan & Sliding-window & 0,6345 & \textbf{0,1507} & \textbf{0,5042} & 0,4626\\
\bottomrule
\end{tabular}
\end{table*}

Hasil tersebut menunjukkan bahwa batas struktur membawa informasi yang diperlukan bersama isi jawaban. Pada manual, parent context mempertahankan judul dan bagian prosedur; pada hukum, child chunk tetap dapat dicari tetapi parent Pasal tersedia ketika konteks dipulihkan. Keunggulan tidak muncul pada semua metrik: sliding-window memiliki MRR sedikit lebih tinggi pada manual dan NDCG sedikit lebih tinggi pada hukum. Hal ini memperlihatkan perbedaan antara menemukan banyak reference dan menempatkan satu hit pertama di posisi awal.

\subsection{Dampak Jenis \indexer{}}

Tabel~\ref{tab:hasil-indexer-paper} memperlihatkan bahwa dense menjadi pilihan terbaik ketika tujuan utama adalah kelengkapan reference. Pada \textit{User Manual}, dense mencapai recall 0,785, precision 0,057, NDCG 0,350, dan MRR 0,314. Nilai tersebut mengungguli sparse dengan selisih recall 0,365. Pertanyaan Indonesia juga mengalami penurunan yang lebih kecil pada dense (recall 0,760) dibandingkan sparse (0,300), yang menunjukkan manfaat representasi multibahasa pada korpus manual berbahasa Inggris.

\begin{table*}[t]
\centering
\caption{Hasil perbandingan \indexer{} pada cutoff $K=30$}
\label{tab:hasil-indexer-paper}
\scriptsize
\begin{tabular}{@{}llrrrr@{}}
\toprule
\textbf{Korpus} & \textbf{Strategi} & \textbf{Recall} & \textbf{Precision} & \textbf{NDCG} & \textbf{MRR}\\
\midrule
User Manual & Structure-aware dense & \textbf{0,785} & \textbf{0,057} & \textbf{0,350} & \textbf{0,314}\\
User Manual & Structure-aware hybrid & 0,740 & 0,048 & 0,333 & 0,231\\
User Manual & Structure-aware sparse & 0,420 & 0,028 & 0,157 & 0,120\\
Peraturan & Legal structure-aware dense & \textbf{0,6854} & \textbf{0,1520} & \textbf{0,5040} & 0,5240\\
Peraturan & Structure-aware hybrid & 0,6412 & 0,1307 & 0,4797 & \textbf{0,5439}\\
Peraturan & Structure-aware sparse & 0,6381 & 0,1253 & 0,4262 & 0,4613\\
\bottomrule
\end{tabular}
\end{table*}

Pada peraturan, dense menghasilkan recall dan NDCG tertinggi, sedangkan hybrid menghasilkan MRR tertinggi. Sinyal lexical membantu menempatkan satu hit pertama pada posisi awal, tetapi tidak selalu menambah seluruh reference yang diperlukan dalam 30 kandidat. Oleh karena itu, pemilihan komponen harus menyatakan tujuan optimasi: dense dipilih untuk coverage, sedangkan hybrid dapat dipertimbangkan ketika hit pertama lebih penting.

\begin{figure}[t]
\centering
\includegraphics[width=\columnwidth]{metrics-at-sweet-spot-indexer.png}
\caption{Ringkasan metrik \indexer{} pada \textit{User Manual} dengan $K=30$}
\label{fig:hasil-indexer}
\end{figure}

\subsection{Perluasan melalui \kg{}}

Eksperimen graf menggunakan 17 relasi berarah dan 24 relasi tak berarah. Seed berasal dari kandidat retrieval langsung, kemudian traversal dilakukan sampai kedalaman dua. Tabel~\ref{tab:hasil-kg-paper} menunjukkan bahwa sliding-window mencapai jangkauan lingkungan penuh, 1,0000, sedangkan recursive dan hybrid mencapai 0,9975. Dense berada pada 0,9556 dan sparse pada 0,9672.

\begin{table}[t]
\centering
\caption{Recall lingkungan \kg{} hukum pada $K=30$}
\label{tab:hasil-kg-paper}
\scriptsize
\begin{tabular}{@{}lr@{}}
\toprule
\textbf{Konfigurasi seed} & \textbf{$R^{\mathrm{KG}}@30$}\\
\midrule
Legal structure-aware dense & 0,9556\\
Recursive & 0,9975\\
Sliding-window & \textbf{1,0000}\\
Structure-aware sparse & 0,9672\\
Structure-aware hybrid & 0,9975\\
\bottomrule
\end{tabular}
\end{table}

Nilai graf yang tinggi tidak berarti bahwa retrieval langsung sudah menemukan seluruh context reference. Graf mengukur apakah dokumen tetangga yang terhubung dapat dijangkau dari seed, sehingga satu seed yang tepat dapat membuka beberapa dokumen lain. Karena itu, \kg{} diposisikan sebagai ekspansi konteks dan pemeriksaan hubungan, bukan pengganti indeks teks.

\begin{figure}[t]
\centering
\includegraphics[width=\columnwidth]{kg-expanded-full-recall-k30.png}
\caption{Recall lingkungan graf untuk konfigurasi hukum pada $K=30$}
\label{fig:hasil-kg}
\end{figure}

\subsection{Dampak Hilir dan Monitoring}

Sebagai pemeriksaan hilir, sepuluh evidence teratas dikirim ke model generatif. Pada \textit{User Manual}, konfigurasi dense menghasilkan correctness 0,788 dan faithfulness 0,958. Pada korpus hukum, legal structure-aware dense menghasilkan correctness 0,690 dan faithfulness 0,972. Angka ini tidak digunakan untuk mengklaim kualitas chatbot secara keseluruhan; angka tersebut hanya menunjukkan bahwa coverage yang lebih baik pada tahap basis pengetahuan dapat diteruskan menjadi evidence yang lebih lengkap.

Monitoring run memberikan konteks untuk membaca hasil tersebut. Setiap baris hasil dapat ditelusuri ke datasource, \profile{}, versi \plugin{}, konfigurasi, ID chunk, dan attempt materialisasi. Dengan demikian, perbedaan skor tidak berasal dari pencampuran indeks lama dan baru. Pencatatan status, jumlah keluaran, penolakan, durasi, dan kesalahan juga memungkinkan proses yang gagal diulang dari artefak terakhir yang valid. Paper ini mengevaluasi fungsi reproducibility tersebut secara kualitatif melalui keberhasilan pengulangan konfigurasi, bukan melalui benchmark latency produksi.

\subsection{Pembahasan}

Temuan utama dapat diringkas menjadi tiga implikasi. Pertama, struktur dokumen adalah bagian dari representasi basis pengetahuan, bukan sekadar metadata. Kedua, dense dan hybrid menjawab tujuan yang berbeda: dense unggul pada coverage, sedangkan hybrid dapat unggul pada hit pertama. Ketiga, \kg{} memberi nilai ketika pertanyaan membutuhkan hubungan lintas dokumen, tetapi metrik graf harus dipisahkan dari metrik relevance teks.

Arsitektur \plugin{} membuat temuan tersebut dapat diuji tanpa menulis ulang orkestrator. Strategi baru cukup mengimplementasikan kontrak yang sama, didaftarkan pada registry, lalu dipilih melalui \profile{}. Keterbatasan arsitektur ini adalah kontrak tidak otomatis menjamin kualitas strategi: validasi skema dan provenance mencegah artefak rusak, tetapi pemilihan konfigurasi tetap memerlukan eksperimen pada domain dan pertanyaan yang relevan.

\subsection{Keterbatasan}

Eksperimen memiliki beberapa batasan. Korpus manual hanya mencakup sembilan produk dan pertanyaannya dibentuk dari konteks acuan, sedangkan korpus hukum hanya mewakili bidang pendidikan dan 25 pertanyaan. OFAT belum mengukur interaksi penuh antara semua kombinasi \chunker{}, \indexer{}, dan parameter traversal. Tidak ada klaim signifikansi statistik karena unit pertanyaan berpasangan dan sebagian reference berbagi sumber. Metrik graf juga mengukur jangkauan lingkungan, bukan validitas hukum atau kecukupan jawaban. Hasil sebaiknya dibaca sebagai bukti untuk konfigurasi yang diuji, bukan sebagai performa universal semua domain.

